%! TeX program = lualatex
% The two comparison-property exercises in lessons/the-definite-integral.tex.
%
%   page 1 -- f(x) <= g(x) over [a,b]
%   page 2 -- m <= f(x) <= M over [a,b], with the two levels y = m and
%             y = M drawn
%
% Both pages use the same window and the same interval [a,b] = [1,5], so the
% pair reads as one idea seen twice.  The axes are ticked only at a and b (and
% at m and M on page 2): the exercises are about comparing net areas, and a
% numbered grid invites counting squares instead.
%
% Every curve crosses the axis inside [a,b] on purpose.  If the curves stayed
% above the axis the comparison would be a statement about plain area and
% would not need the net area reading the exercises ask for; as drawn, each
% integral is a positive piece minus a negative one.
%
% Equal 1.2cm units on both axes, as elsewhere in these notes.  Include the
% pages at natural size -- no width= -- or the units stop matching.
\documentclass[tikz, crop]{standalone}

\input{../typography.tex.preamble}
\input{../tikz.tex.preamble}

\pgfplotsset{comparison panel/.style={
    basic curve,
    axis equal image=false,
    set layers,
    x=1.2cm, y=1.2cm,
    xmin=-0.5, xmax=6.6,
    xlabel={\(x\)}, ylabel={},
    grid=none,
    clip=false,
  }}

\begin{document}

% ---------------------------------------------------------------------
% Page 1 -- f(x) <= g(x).
%
% f is g shifted down by 0.8 + 0.3 sin(1.7x).  Building f from g this way
% guarantees the inequality everywhere -- the gap runs between 0.5 and 1.1 --
% while the second sine keeps the two curves from being parallel copies, which
% would suggest the difference has to be constant.
%
% Over [1,5]: g runs from 1.42 down to -0.87 and crosses at x = 3.84, f runs
% from 0.32 down to -1.91 and crosses at x = 3.31.
\begin{tikzpicture}
  \begin{axis}[
    comparison panel,
    ymin=-2.3, ymax=2.1,
    xtick={1,5}, xticklabels={\(a\),\(b\)}, ytick=\empty,
    ]
    \draw[dashed, gray] (axis cs:1,-2.3) -- (axis cs:1,2.1);
    \draw[dashed, gray] (axis cs:5,-2.3) -- (axis cs:5,2.1);
    \addplot[very thick, domain=1:5] {1.3*sin(0.9*x) + 0.4};
    \addplot[very thick, domain=1:5]
      {1.3*sin(0.9*x) + 0.4 - (0.8 + 0.3*sin(1.7*x))};
    \node[anchor=west] at (axis cs:5.1,-0.87) {\(y = g(x)\)};
    \node[anchor=west] at (axis cs:5.1,-1.91) {\(y = f(x)\)};
  \end{axis}
\end{tikzpicture}

% ---------------------------------------------------------------------
% Page 2 -- m <= f(x) <= M.
%
% f(x) = 1.2 sin(1.05x) + 0.3 runs between -0.90 and 1.50 over [1,5] and
% crosses the axis at x = 3.24, so the bounds m = -1 and M = 2 hold with a
% little room either side: m and M are bounds, not the extreme values.
%
% m is negative so that the lower level sits below the axis: m(b-a) = -4 is a
% negative net area, which is the whole point of reading the inequality that
% way.  The upper level gives M(b-a) = 8.
%
% Only the two horizontal levels are drawn, not a rectangle: the regions
% whose net areas are M(b-a) and m(b-a) are what the reader has to find, and
% outlining them gives that away.  Each level runs from the y axis, where its
% label sits, to b; a and b are marked by full-height verticals, as on page 1.
\begin{tikzpicture}
  \begin{axis}[
    comparison panel,
    ymin=-1.8, ymax=2.6,
    xtick={1,5}, xticklabels={\(a\),\(b\)},
    ytick={-1,2}, yticklabels={\(m\),\(M\)},
    ]
    \draw[dashed, gray] (axis cs:1,-1.8) -- (axis cs:1,2.6);
    \draw[dashed, gray] (axis cs:5,-1.8) -- (axis cs:5,2.6);
    \draw[dashed, gray] (axis cs:0,2)  -- (axis cs:5,2);
    \draw[dashed, gray] (axis cs:0,-1) -- (axis cs:5,-1);
    \addplot[very thick, domain=1:5] {1.2*sin(1.05*x) + 0.3};
    \node[anchor=west] at (axis cs:5.1,-0.73) {\(y = f(x)\)};
  \end{axis}
\end{tikzpicture}

\end{document}
